\section{Interstate 2.0 Integration}
\label{sec:i2}

\subsection{The Port Connector as the I2.0 On-Ramp}

Interstate 2.0 managed freight lanes on trunk T1 corridors are designed to operate at high speed (65--70 mph) with geometric consistency (12-foot lanes, 10-foot shoulders, limited access interchanges). The value proposition is that a truck entering the managed lane at one relay terminal can travel to the next relay terminal with predictable transit times and no mixing with slower traffic. The managed-lane system is, in effect, a premium freight highway that charges for the certainty premium it provides.

A port connector that enters the managed-lane network at V/C 1.5 breaks this value proposition immediately. A truck exiting the Port of Long Beach at 6:00~am in the I2.0 managed-lane network faces a 45-minute crawl on the I-710 before it can enter the managed lane at the I-710/I-405 junction. The managed-lane reliability that the shipper paid for begins at the I-405 interchange; the 45-minute port connector delay is invisible to the managed-lane performance metric and cannot be SLA-committed. Port connector congestion leaks into the managed-lane performance envelope.

The I2.0 system design should therefore treat the port connector as the first segment of the managed-lane corridor for performance guarantee purposes. This has two operational implications:

\begin{enumerate}
  \item \textbf{V/C standard consistency}: The I2.0 managed-lane design target is V/C $\leq 0.85$ at the 95th-percentile peak hour. Port connectors serving I2.0-designated corridors should be required to meet the same V/C standard. A port connector at V/C 1.5 is not a T1 bottleneck---it is an I2.0 bottleneck that happens to be 11 miles long.

  \item \textbf{Drayage optimization integration}: The relay hub load-matching platform (Section~\ref{sec:anatomy} of the companion paper \citep{ROUTE_C3}) can integrate port gate scheduling data to optimize drayage sequencing on the port connector. If the load-matching algorithm knows that 400 trucks will arrive from the port between 9:00--10:00~am, it can sequence the gate releases to spread departures across a 2-hour window---reducing the peak-hour demand on the connector from 14,000 trucks/hour to 8,000--9,000 trucks/hour without adding any capacity. This operational smoothing can reduce V/C on the I-710 from 1.52 to approximately 1.15 at no infrastructure cost, as a short-term measure while physical improvements proceed through the NEPA process.
\end{enumerate}

\subsection{Drayage Optimization as Demand Management}

Port drayage involves two classes of truck movement: live-load moves (truck arrives at port, loads a container, departs immediately) and street-turn moves (truck arrives with an empty container, exchanges for a loaded outbound container at the terminal). The latter requires less dwell time and generates half the number of connector trips per container move.

The Georgia Ports Authority, Port of Long Beach, and Port of New York/New Jersey all operate appointment systems that schedule truck gate arrivals in 30-minute windows. These appointment systems currently optimize for terminal throughput (reducing yard congestion), not connector throughput. Integrating appointment scheduling with real-time connector V/C monitoring would allow the terminal operator to smooth connector demand: when the V/C exceeds 1.1, the appointment system holds the next wave of outbound gate releases for 15--20 minutes, allowing the connector to clear.

This demand management approach reduces peak connector V/C by approximately 20--25\% in simulation \citep{POLA_gate2023} without requiring physical infrastructure changes. It is a prerequisite for I2.0 integration because it provides the real-time demand signal that the managed-lane control system needs to manage entry onto the T1 corridor.

\subsection{The I2.0 Port Connector Design Standard}

An I2.0-compatible port connector must meet the following minimum standards, which we propose as the ``port connector'' classification criteria:

\begin{itemize}
  \item \textbf{Geometric standard}: Limited-access freeway standard (no at-grade intersections) for the full length of the connector; 12-foot lanes; 10-foot right shoulder; interchange access only (no continuous right-turn movements).
  \item \textbf{V/C standard}: V/C $\leq 0.85$ at the 95th-percentile peak hour, measured using the adjusted port peak-hour methodology (Equation~\ref{eq:vc} above), not 24-hour AADT.
  \item \textbf{Resilience standard}: At least one viable alternate connector route capable of handling 60\% of normal truck volumes in the event of primary connector closure; alternate route identified and signed.
  \item \textbf{Data standard}: Real-time V/C monitoring with 5-minute update frequency, integrated with the port gate appointment system and the I2.0 managed-lane control system.
  \item \textbf{NHS designation}: The connector must be on the National Highway System to qualify for NHPP funding and I2.0 program inclusion.
\end{itemize}

Currently, only two of the ten connectors analyzed in this paper meet all five criteria: the I-10/I-610 connector at Houston (meets geometric, NHS, and data standards but not V/C at peak) and the I-64 Norfolk connector (meets geometric and NHS standards but not V/C or resilience). The other eight connectors require significant investment to reach I2.0 compatibility.

\subsection{Integration with Relay Hub Siting}

The relay hub siting analysis in \citet{ROUTE_C1} identifies optimal hub locations based on truck AADT, flow imbalance, and HOS constraints. Port connectors introduce an additional siting criterion: hubs near major ports should be positioned to capture port-origin freight that would otherwise deadhead on the connector.

A relay hub at Carson, CA (between the I-710 terminus and the I-405/I-710 interchange) would capture trucks exiting the Port of Long Beach and offer immediate load-matching for the inland haul. This hub location is not optimal under the standard flow-balance criterion (the LA metro is relatively balanced in freight flows), but it is highly valuable for reducing connector congestion: a truck that matches a load at the Carson hub before entering the I-710 peak-flow period can delay its departure by 2--3 hours, naturally smoothing the connector peak demand. The hub's load-matching function thus acts as demand management for the connector---a synergy not captured in either the relay hub analysis alone or the port connector analysis alone.
